\documentclass[tikz,border=14pt]{standalone}
\usepackage{amsmath}
\usetikzlibrary{calc}
\definecolor{bg}{HTML}{171D2B}
\definecolor{aksen}{HTML}{F7B538}
\begin{document}
\pagecolor{bg}
\begin{tikzpicture}[line width=1.2pt, text=white, scale=1.0]
  % garis l1 dan l2 (sejajar), transversal m
  \coordinate (L1a) at (0,0);
  \coordinate (L1b) at (9,0);
  \coordinate (L2a) at (-0.8,-4);
  \coordinate (L2b) at (8.2,-4);
  \coordinate (Ta) at (1.0,1.8);
  \coordinate (Tb) at (7.0,-5.8);

  \draw[color=aksen] (L1a) node[left]{\large $l_1$} -- (L1b);
  \draw[color=aksen] (L2a) node[left]{\large $l_2$} -- (L2b);
  \draw[white] (Ta) node[above]{\large $m$} -- (Tb);

  % titik potong
  \coordinate (P) at (3,0);
  \coordinate (Q) at (5,-4);

  % sudut sehadap di atas garis: 1 (di P, atas-kanan) dan 5 (di Q, atas-kanan)
  \draw[aksen] ($(P)+(0.9,0)$) arc[start angle=0, end angle=60, radius=0.9];
  \node[aksen] at ($(P)+(40:1.35)$) {\large $\alpha$};
  \draw[aksen] ($(Q)+(0.9,0)$) arc[start angle=0, end angle=60, radius=0.9];
  \node[aksen] at ($(Q)+(40:1.35)$) {\large $\alpha$};

  % sudut bersudut berhadapan: di Q, atas-kiri (dalam) 120 derajat
  \draw[white] ($(Q)+(-0.7,0)$) arc[start angle=180, end angle=120, radius=0.7];
  \node[white] at ($(Q)+(150:1.45)$) {\large $180^{\circ}-\alpha$};

  % sudut dalam sepihak di P: kiri-bawah (di bawah l1, kiri transversal)
  \draw[white] ($(P)+(-0.7,0)$) arc[start angle=180, end angle=240, radius=0.7];
  \node[white] at ($(P)+(215:1.45)$) {\large $180^{\circ}-\alpha$};

  % tanda sejajar (panah) pada l1 dan l2
  \draw[aksen] ($(L1a)!0.5!(L1b)+( -0.1,0.12)$) -- ($(L1a)!0.5!(L1b)+(0.1,0)$) -- ($(L1a)!0.5!(L1b)+(-0.1,-0.12)$);
  \draw[aksen] ($(L2a)!0.5!(L2b)+(-0.1,0.12)$) -- ($(L2a)!0.5!(L2b)+(0.1,0)$) -- ($(L2a)!0.5!(L2b)+(-0.1,-0.12)$);
\end{tikzpicture}
\end{document}
